\SourceLeaf{396}%
and
\[
\tfrac{c+d+e}{3}-\tfrac13\BareRoot{(cc+dd+ee+cd+ce+de)}
\]

The radical expression \(\BareRoot{(cc+dd+ee+cd+ce+de)}\) can also be written in the form
\[
\BareRoot{[(d-c)^2+(d-c)(e-d)+(e-d)^2]}
\]
It is therefore the third side of a triangle with two sides \(d-c\) and \(e-d\) and included angle \(=120^0\). If we construct on \(CD\) an equilateral triangle with apex \(F\), then \(EF\) equals the radical expression, so the two values of \(r\) can readily be constructed. It is easy to show that one of these values must lie between \(c\) and \(d\) and the other between \(d\) and \(e\), and that only the former actually gives the centre of the largest ellipse. For the other value,
\[
\gamma+\delta r+\varepsilon rr+\zeta r^3
\]
will attain a minimum instead of a maximum, or rather will take its greatest \textit{negative} value, to which only an imaginary value of \(\alpha\beta\) can correspond. It is easy to see that this must refer to a hyperbola.

Once the centre of the required ellipse has been found, determining the remaining unknown quantities presents no difficulty. From \(\theta\) and \(r\) we find \(t\). From \(t\) and \(u\) we then find \(\alpha\) and \(\beta\), and from one or more of the equations above we find \(\psi\). This determines both the dimensions and the position of the ellipse completely.

I must also point out that the problem solved here differs slightly from the one recently posed in the \textit{Monatl. Corresp.} There are cases in which the largest ellipse that can be drawn \textit{inside} a quadrilateral does not touch one of its four sides. A closer examination of those cases is beyond my purpose here.

\par\medskip\Needspace{4\baselineskip}
\pdfbookmark[2]{Direct proof of the theorem above concerning quadrilaterals.}{ellipse-directer-beweis}
\begin{center}\textit{Direct proof of the theorem above concerning quadrilaterals.}\end{center}
\hspace*{\parindent}Let \(A,B,C,D\) be the four vertices of the quadrilateral. Let \(E\) be the intersection of \(AB\) and \(DC\), let \(F\) be the intersection of \(BC\) and \(AD\), and let \(G,H\) and \(I\)
